\documentclass[9pt]{beamer}

\usecolortheme[RGB={205,173,0}]{structure} 
\usefonttheme{serif} 

\usepackage{amssymb}
\usepackage[all]{xy}

%Typesetting conventions
\def\K3{\mathrm K3}
\def\CY#1{\mathrm{CY}_{#1}}
%\def\SU2{\mathrm{SU}(2)}
\def\SU#1{SU({#1})}
%\def\U1{\mathrm U(1)}
\def\U#1{U({#1})}
\def\SO#1{SO({#1})}
\def\O#1{O({#1})}
%multiplet and dimension
\def\mult #1#2{\mathrm{#1}[\![ {#2}]\!]}
%normal ordering
\def\normal#1{\,\!:\! #1 \!:\! \,}
%doublebar
\def\double #1{#1{\hbox{\kern-2pt $#1$}}}
%underline
\def\un#1{\underline #1}


%%%%%%%%%%%%%%%%%%%%%%%%%%%%%%%%%%%%%%%%%%%%%%%%
%%%%%%%%%%%%%%%%  Typesetting conventions  %%%%%%%%%%%%%%%%%%%

%doublebar
\def\double #1{#1{\hbox{\kern-2pt $#1$}}}
%Bras and kets
\def\ket#1{\left| #1\right>}
\def\bra#1{\left< #1\right|}
\def\braket#1#2{\left< #1\right|\left.#2\right>}
%For twistors
\def\kra#1{\left[ #1\right|}
\def\bet#1{\left| #1\right]}

%%%%%%%%%%%%%%%%%%%%%%%%%%%%%%%%%%%%%%%%%%%%%%%%
%%%%%%%%%%%%  Jim Gates and Marc Grisaru MACROS   %%%%%%%%%%%%%%%%
\def\pp{{\mathchoice
            %{general format
               %[w] = length of horizontal bars
               %[t] = thickness of the lines
               %[h] = length of the vertical line
               %[s] = spacing around the symbol
              %
              %\kern [s] pt%
              %\raise 1pt
              %\vbox{\hrule width [w] pt height [t] pt depth0pt
              %      \kern -([h]/3) pt
              %      \hbox{\kern ([w]-[t])/2 pt
              %            \vrule width [t] pt height [h] pt depth0pt
              %            }
              %      \kern -([h]/3) pt
              %      \hrule width [w] pt height [t] pt depth0pt}%
              %      \kern [s] pt
          {%displaystyle
              \kern 1pt%
              \raise 1pt
              \vbox{\hrule width5pt height0.4pt depth0pt
                    \kern -2pt
                    \hbox{\kern 2.3pt
                          \vrule width0.4pt height6pt depth0pt
                          }
                    \kern -2pt
                    \hrule width5pt height0.4pt depth0pt}%
                    \kern 1pt
           }
            {%textstyle
              \kern 1pt%
              \raise 1pt
              \vbox{\hrule width4.3pt height0.4pt depth0pt
                    \kern -1.8pt
                    \hbox{\kern 1.95pt
                          \vrule width0.4pt height5.4pt depth0pt
                          }
                    \kern -1.8pt
                    \hrule width4.3pt height0.4pt depth0pt}%
                    \kern 1pt
            }
            {%scriptstyle
              \kern 0.5pt%
              \raise 1pt
              \vbox{\hrule width4.0pt height0.3pt depth0pt
                    \kern -1.9pt  %[e]=0.15pt
                    \hbox{\kern 1.85pt
                          \vrule width0.3pt height5.7pt depth0pt
                          }
                    \kern -1.9pt
                    \hrule width4.0pt height0.3pt depth0pt}%
                    \kern 0.5pt
            }
            {%scriptscriptstyle
              \kern 0.5pt%
              \raise 1pt
              \vbox{\hrule width3.6pt height0.3pt depth0pt
                    \kern -1.5pt
                    \hbox{\kern 1.65pt
                          \vrule width0.3pt height4.5pt depth0pt
                          }
                    \kern -1.5pt
                    \hrule width3.6pt height0.3pt depth0pt}%
                    \kern 0.5pt%}
            }
        }}

\def\mm{{\mathchoice
                      %{general format %[w] = length of bars
                                       %[t] = thickness of bars
                                       %[g] = gap between bars
                                       %[s] = space around symbol
   %[w], [t], [s], [h]=3([g]) are taken from corresponding definitions of \pp
   %
                      %       \kern [s] pt
               %\raise 1pt    \vbox{\hrule width [w] pt height [t] pt depth0pt
               %                   \kern [g] pt
               %                   \hrule width [w] pt height[t] depth0pt}
               %              \kern [s] pt}
                  %
                       {%displaystyle
                             \kern 1pt
               \raise 1pt    \vbox{\hrule width5pt height0.4pt depth0pt
                                  \kern 2pt
                                  \hrule width5pt height0.4pt depth0pt}
                             \kern 1pt}
                       {%textstyle
                            \kern 1pt
               \raise 1pt \vbox{\hrule width4.3pt height0.4pt depth0pt
                                  \kern 1.8pt
                                  \hrule width4.3pt height0.4pt depth0pt}
                             \kern 1pt}
                       {%scriptstyle
                            \kern 0.5pt
               \raise 1pt
                            \vbox{\hrule width4.0pt height0.3pt depth0pt
                                  \kern 1.9pt
                                  \hrule width4.0pt height0.3pt depth0pt}
                            \kern 1pt}
                       {%scriptscriptstyle
                           \kern 0.5pt
             \raise 1pt  \vbox{\hrule width3.6pt height0.3pt depth0pt
                                  \kern 1.5pt
                                  \hrule width3.6pt height0.3pt depth0pt}
                           \kern 0.5pt}
                       }}

\def\ad{{\kern0.5pt
                   \alpha \kern-5.05pt
\raise5.8pt\hbox{$\textstyle.$}\kern 0.5pt}}

\def\bd{{\kern0.5pt
                   \beta \kern-5.05pt \raise5.8pt\hbox{$\textstyle.$}\kern 0.5pt}}

\def\qd{{\kern0.5pt
                   q \kern-5.05pt \raise5.8pt\hbox{$\textstyle.$}\kern 0.5pt}}
\def\Dot#1{{\kern0.5pt
     {#1} \kern-5.05pt \raise5.8pt\hbox{$\textstyle.$}\kern 0.5pt}}
%%%%%%%%%%%%%%%%%%%%%%%%%%%%%%%%%%%%%%%%%%%%%%%%%%%
%%%%%%%%%%%%%%%%%%%%%%%%%%%%%%%%%%%%%%%%%%%%%%%%%%%


\usepackage{beamerthemesplit}

\title{Superstrings in Flux Backgrounds} %\\(Proyecto n\'umero 11100425)}
\author{William D. Linch III}
\date{21 diciembre, 2011.}

\begin{document}

\frame{\titlepage}

%\section[Outline]{}
%\frame{\tableofcontents}

%%%%%%%%%%%%%%%%%%%%%%%%%%%%%%%%%%%%%%%%%%%%%%%%%%%
\section{Superstrings in Flux Backgrounds}
\frame{
\frametitle{F\'isca Fundamental Actual y en el Futuro Cercano}
Hay mucha emoci\'on reciente sobre la f\'isica de part\'iculas fundamentales%\pause :
\begin{itemize}
\item T2K: >Oscilaci\'on de neutrinos $\nu_\mu \to \nu_e$? ($\sim2.5\sigma$) 
\item OPERA: >>>Propagaci\'on superlum\'inico de neutrinos??? ($\sim6\sigma$)
\item LHCb: >Exceso de violaci\'on de CP en $D^0\to K^+K^-$ y $D^0\to \pi^+\pi^-$? ($\sim3.5\sigma$)
\item LHC: >Exceso de $pp\to \ell^-\bar\nu\ell^+\nu$ consistente con Higgs de $m_h\sim 125$ GeV? (ATLAS:$\sim2.6\sigma$ y CMS: $\sim1.9\sigma$)
\end{itemize}
%\pause  
Confirmaci\'on y violaci\'on del Modelo Est\'andar simulaneamente!\\
$\Rightarrow$ F\'isica m\'as all\'a del Modelo Est\'andar%\pause :
\begin{itemize}
\item >Dimensiones adicionales?
\item >Gran Unificaci\'on de electromagnetismo, la fuerza d\'ebil, y la fuerza fuerte? (GUT: $SU(5)$, $SO(10)$, $E_6$, \dots)
\item >Supersimetr\'ia? (SUSY/SUGRA)
\item >...?
\end{itemize}
}

\frame{
\frametitle{El Espacio de Teor\'ias Fundamentales}
\begin{displaymath}
\xymatrix@=10pt{ 
{\begin{array}{c}\textrm{Grav.}\\ \textrm{Cu\'ant.}\end{array}}&& \ar[ll]{\begin{array}{c}\textrm{Camp.}\\ \textrm{Cu\'ant.}\end{array}}& \\
&&&{\begin{array}{c}\textrm{Mec.}\\ \textrm{Cu\'ant.}\end{array}}\ar[ul] \\
{\begin{array}{c}\textrm{Rel.}\\ \textrm{Gen.}\end{array}}\ar[uu] &  & 
	{\begin{array}{c}\textrm{Rel.}\\ \textrm{Esp.}\end{array}}\ar[ll]\ar[uu]& \\
&{\begin{array}{c}\textrm{Grav.}\\ \textrm{Cl\'as.}\end{array}}\ar[ul]&&
	\ar[ll]_{G}{\begin{array}{c}\textrm{Mec.}\\ \textrm{Cl\'as.}\end{array}} \ar[ul]_{1/c}\ar[uu]_\hbar
}
\end{displaymath}
}

\frame{
\frametitle{El Espacio de Teor\'ias Fundamentales}
\begin{displaymath}
\xymatrix@=10pt{ 
{\begin{array}{c}\textrm{Grav.}\\ \textrm{Cu\'ant.}\end{array}}&& \ar[ll]{\color{red}\begin{array}{c}\textrm{Camp.}\\ \textrm{Cu\'ant.}\end{array}}& \\
&{\color{red}\begin{array}{c}\textrm{Mod.}\\ \textrm{Est.}\end{array}} \ar@{=>}[ur]\ar@{:>}[dl]&&{\begin{array}{c}\textrm{Mec.}\\ \textrm{Cu\'ant.}\end{array}}\ar[ul] \\
{\color{red}\begin{array}{c}\textrm{Rel.}\\ \textrm{Gen.}\end{array}}\ar[uu] &  & 
	{\begin{array}{c}\textrm{Rel.}\\ \textrm{Esp.}\end{array}}\ar[ll]\ar[uu]& \\
&{\begin{array}{c}\textrm{Grav.}\\ \textrm{Cl\'as.}\end{array}}\ar[ul]&&
	\ar[ll]_{G}{\begin{array}{c}\textrm{Mec.}\\ \textrm{Cl\'as.}\end{array}} \ar[ul]_{1/c}\ar[uu]_\hbar
}
\end{displaymath}
}

\frame{
\frametitle{Supercuerdas}
Necesidad de ``cuerdas'':%\pause 
\begin{itemize}
\item Cuantizaci\'on de part\'iculas relativ\'isticas $\Rightarrow$ Teor\'ia Cu\'antica de Campos = el c\'alculo del Modelo Est\'andar ($\hbar$ \& $c$)
%%\pause  
\item Cuantizaci\'on de cuerdas relativ\'isticas $\Rightarrow$ Teor\'ia de todas part\'iculas y fuerzas ($\hbar$, $c$, \& $G$)
%%\pause  
\item Formulaci\'on sencillo $\Rightarrow$ $D=10$ \& {SUSY}%\pause :
%{\bf I}, {\bf II}{\sc a}, {\bf II}{\sc b}, {\bf H}{\sc e}, {\bf H}{\sc o}
\end{itemize}
%~\\
\begin{displaymath}
\xymatrix@=10pt{ 
&\ar@{-}@/^1pc/[dl]{\begin{array}{c}11\textrm{\sc d}\\ \textrm{{\sc sugra}}\end{array}} \ar@{-}@/_1pc/[dr]& \\
\mathbf {II}\textrm{\sc a}\ar@{-}@/^1pc/[dd]&&\mathbf {H}\textrm{\sc e}\ar@{-}@/_1pc/[dd]\\
&\mathbf M&\\
\mathbf {II}\textrm{\sc b}&&\mathbf {H}\textrm{\sc o}\\
&\ar@{-}@/_1pc/[ul]{\mathbf I}\ar@{-}@/^1pc/[ur]&
}
\end{displaymath}
}

\frame{
\frametitle{Supercuerdas}
Formulaciones de supercuerdas:
\begin{itemize}
\item Ramond-Neveu-Schwarz 
\item Green-Schwarz 
\item H\'ibrida $=$ (Ramond-Neveu-Schwarz) $\oplus$ (Green-Schwarz)
\item Covariante $\supset$ Green-Schwarz
\end{itemize}
%\pause  
Implementaci\'on en c\'alculos fen\'omenologicos:
\begin{itemize}
\item $D<10$ $\Rightarrow$ ``compactificaci\'on'' y ``holograf\'ia''
\item Fondos con flujos no-triviales
\item $\Rightarrow$ Teor\'ia de cuerdas = Herramienta computacional muy poderoso pero tambi\'en {\em muy complicado}
\end{itemize}
}

%\frame{
%\frametitle{Supercuerdas}
%\begin{displaymath}
%\xymatrix{ 
%???&&Covariante \\
%&Hibrida\ar[u] &\\
%RNS\ar[ur] &  & GS\ar[ul]
%}
%\end{displaymath}
%}


\frame{
\frametitle{Proyecto propuesto (11100425)}
Fase I
\begin{itemize}
\item Formulaci\'on de la supercuerda covariante (tipo {\bf I}, {\bf H}{\sc e}, {\bf H}{\sc o}) en un fondo Calabi-Yau
\item Establecer relaci\'on con forumulaci\'on ``h\'ibrida'' en $D=4$ y $D=6$
\end{itemize}
%%\pause  
Fase II
\begin{itemize}
\item Investigaci\'on de soluciones de problemas t\'ecnicos en formulaciones no covariantes
\item Investigar posibilidad de simetr\'ia especular en las cuerdas {\bf H}{\sc e} y {\bf H}{\sc o}
\end{itemize}
%%\pause  
Fase III
\begin{itemize}
\item Inclusi\'on de flujos en fondos tipo {\bf II}{\sc a} y {\bf II}{\sc b} para fenomenolog\'ia
\item Descripci\'on de part\'iculas en fondos hologr\'aficos $AdS\times X$
\item Uso de resultados para aclarar estructura de superespacios y agujeros negros
\end{itemize}
}

\frame{
\frametitle{Fase I: Resultados}
\begin{itemize}
\item Compactificaciones a $D=4,6,$ y 8 en ``orbifoldes'' con ``holonom\'ia especial''
	\begin{itemize}
	\item {\em Compactification of the Heterotic Pure Spinor Superstring II}, JHEP {\bf 1110}, 098 (2011) [arXiv:1108.3555 [hep-th]] junto con Chand\'ia ({\sc uai}) y Vallilo ({\sc unab}).
	\end{itemize}
\item Compactificacion a $D=6$ en fondo de K3
\item Relaci\'on de la cuerda en K3 con descripci\'on h\'ibrida 
	\begin{itemize}
	\item {\em The Covariant Superstring on K3}, [arXiv:1109.3200 [hep-th]] con Chand\'ia ({\sc uai}) y Vallilo ({\sc unab}).
	\end{itemize}
\item Desarrollo de SUGRA en el superespacio proyectivo/arm\'onico con $D=6$ y $N=(1,0)$
	\begin{itemize}
	\item {\em Super-Weyl invariance in 6D supergravity}, (en preparaci\'on) con G. Tartaglino-Mazzucchelli ({Uppsala}, Swecia).
	\end{itemize}
\end{itemize}
}

\frame{
\frametitle{Fase I: Otros}
\begin{itemize}
\item Presentaci\'on: {\em Compactification of the pure spinor superstring}, 5th {\sc ictp-capes} Latin-American String School (diciembre, 2010), {\sc ift-unesp}, Sa\~o Paulo, Brazil.
\item Visita de colaboraci\'on a Dr Gabriele Tartaglino- Mazzucchelli en Uppsala University (28 julio|11 agosto, 2011), Uppsala, Suecia.
\item Inicio trabajo con alumno de mag\'ister (C\'esar Patricio Arias Huerta) en geometr\'ia de superespacios de seis dimensiones con supersimetr\'ia m\'inimo.
\item Entrevistas sobre neutrinos superlum\'inicos/dimensiones adicionales (23/9/2011 por TVN) y descubrimiento del Higgs/supersimetr\'ia (14/12/2011 en La Tercera).
\end{itemize}
}


%%%%%%%%%%%%%%%%%%%%%%%%%%%%%%%%%%%%%%%%%%%%%%%%
%\frame{
%\frametitle{Fin}
%\begin{centering}
%{\Huge Thank you!}
%\end{centering}
%}

%%%%%%%%%%%%%%%%%%%%%%%%%%%%%%%%%%%%%%%%%%%%%%%%
%\begin{thebibliography}{99}
%
%\footnotesize{
%  
%%\cite{Chandia:2011su}
%\bibitem{Chandia:2011su}
%  O.~Chandia, W.~D.~Linch, III, B.~C.~Vallilo,
%  ``The Covariant Superstring on K3,''
%  [arXiv:1109.3200 [hep-th]].
%
%}%end footnotesize
%
%\end{thebibliography}


\end{document}